\documentclass[11pt,a4paper]{article}

% ── Engine: pdfLaTeX ─────────────────────────────────────────────────────────
\usepackage[utf8]{inputenc}
\usepackage[T1]{fontenc}
\usepackage{lmodern}
\usepackage[margin=2.3cm, headheight=15pt]{geometry}
\usepackage{hyperref}
\usepackage{booktabs}
\usepackage{longtable}
\usepackage{array}
\usepackage{xcolor}
\usepackage{titlesec}
\usepackage{enumitem}
\usepackage{microtype}
\usepackage{parskip}
\usepackage{fancyhdr}
\usepackage{newunicodechar}
\usepackage{amssymb}
\usepackage{tcolorbox}

\hypersetup{colorlinks=true, linkcolor=accent, urlcolor=accent}

% Absorb long unbreakable \texttt{} tokens into interword space instead of
% letting them run into the right margin.
\emergencystretch=3em

% ── Indian Rupee sign (drawn, pdfLaTeX-compatible) ───────────────────────────
\newcommand{\Rs}{%
  \mbox{%
    \sffamily\bfseries
    \setbox0=\hbox{R}%
    R\kern-\wd0
    \raisebox{0.42em}{\rule{0.55em}{0.07em}}%
    \kern-\wd0
    \raisebox{0.25em}{\rule{0.55em}{0.07em}}%
    \kern\wd0
  }%
}
\newunicodechar{₹}{\texorpdfstring{\Rs}{Rs.}}

% ── Colours ──────────────────────────────────────────────────────────────────
\definecolor{accent}{RGB}{30, 80, 160}
\definecolor{mgreen}{RGB}{34, 139, 34}
\definecolor{mgrey}{RGB}{90, 90, 90}
\definecolor{boxbg}{RGB}{238, 243, 250}

% ── Section formatting ───────────────────────────────────────────────────────
\titleformat{\section}{\large\bfseries\color{accent}}{\thesection.}{0.5em}{}[\titlerule]
\titleformat{\subsection}{\normalsize\bfseries}{\thesubsection}{0.6em}{}

% ── Header / footer ──────────────────────────────────────────────────────────
\pagestyle{fancy}
\fancyhf{}
\rhead{\textcolor{mgrey}{\small TestFlow --- Revenue Architecture}}
\lhead{\textcolor{mgrey}{\small 2026-05-30}}
\rfoot{\textcolor{mgrey}{\small Page \thepage}}
\renewcommand{\headrulewidth}{0.4pt}

% ── Column types ─────────────────────────────────────────────────────────────
\newcolumntype{L}[1]{>{\raggedright\arraybackslash}p{#1}}
\newcolumntype{R}[1]{>{\raggedleft\arraybackslash}p{#1}}
\newcolumntype{C}[1]{>{\centering\arraybackslash}p{#1}}

% ── Helpers ──────────────────────────────────────────────────────────────────
\newcommand{\tild}{\raisebox{0.17ex}{$\scriptstyle\sim$}}
\newcommand{\plan}[1]{\texttt{\textcolor{accent}{#1}}}
\newcommand{\flag}[1]{\texttt{#1}}
\newcommand{\yes}{\textcolor{mgreen}{\checkmark}}
\newcommand{\no}{\textcolor{mgrey}{--}}

\newtcolorbox{keybox}{colback=boxbg, colframe=accent, boxrule=0.6pt,
  left=6pt, right=6pt, top=5pt, bottom=5pt, arc=2pt}

% ─────────────────────────────────────────────────────────────────────────────
\begin{document}
% ─────────────────────────────────────────────────────────────────────────────

\begin{center}
  {\LARGE\bfseries TestFlow Revenue Architecture}\\[0.4em]
  {\large\color{accent} Subscription Plans \& Revenue Models}\\[1.2em]
  \begin{tabular}{ll}
    \textbf{Date}    & 2026-05-30 \\
    \textbf{Author}  & Parakh Sharma (solo founder) \\
    \textbf{Status}  & Draft for review \\
    \textbf{Related} & \texttt{PRICING.md}, \texttt{SALES\_TARGETS.md}, financial-plan \texttt{.tex} \\
  \end{tabular}
\end{center}
\vspace{0.8em}
\hrule
\vspace{1.2em}

% ── Overview ─────────────────────────────────────────────────────────────────
\section{Why this document}

\texttt{PRICING.md} commits to one sound idea: a \textbf{flat per-workspace annual
fee}, no per-seat metering. Keep that as the \emph{spine}. But a single flat fee
leaves money --- and customers --- on the table at both ends of the market: the
seasonal two-substation contractor who will never sign an annual contract, and
the utility that wants a framework license for 40 sites.

This document layers \textbf{seven revenue models} onto that spine so every buyer
has a way in and every customer has a way to spend more over time. It is grounded
in what the database already supports: \texttt{subscriptions.plan\_id} (free-text
SKU), \texttt{subscriptions.seat\_count}, \texttt{companies.features} (JSONB
feature flags read by \texttt{has\_feature()} / \texttt{useFeature()}), and the
Razorpay webhook (\texttt{upsert\_subscription}).

\begin{keybox}
\textbf{Design principle --- Land, Expand, Monetize the edges.}
The annual subscription is the \emph{land}. Metered add-ons and seat-free
capacity tiers are the \emph{expand}. Consumption packs, outcome pricing,
white-label, and utility framework deals \emph{monetize the edges} of the market
the flat fee cannot reach.
\end{keybox}

% ── The revenue stack ────────────────────────────────────────────────────────
\section{The revenue stack at a glance}

\begin{center}
\begin{tabular}{@{}C{0.5cm} L{3.9cm} L{5.7cm} L{3.5cm}@{}}
  \toprule
  \textbf{\#} & \textbf{Revenue model} & \textbf{What the customer buys} & \textbf{Recurring?} \\
  \midrule
  1 & Core subscription ladder   & Annual workspace access \& tier features      & Yes (annual) \\
  2 & Hybrid metered add-ons     & Usage beyond fair-use (AI, storage, comms)    & Yes (variable) \\
  3 & Consumption / project packs& Prepaid ``project credits'', no contract      & No (repeat) \\
  4 & Outcome / value pricing    & Per signed handover / commissioned bay        & Yes (variable) \\
  5 & Channel \& white-label     & OEM/EPC resells to subcontractors             & Yes (wholesale) \\
  6 & Utility framework license  & Multi-site, multi-year, on-prem option        & Yes (multi-yr) \\
  7 & Services \& data products  & Onboarding, migration, benchmarks, training   & Mixed \\
  \bottomrule
\end{tabular}
\end{center}

\noindent A customer typically starts in one model and accretes others. A mid-size
contractor lands on \plan{team\_annual} (1), buys AI-report overages (2), takes a
white-label add-on (7), and three years later signs a utility framework (6) when
they win a PowerGrid package.

% ── Model 1 ──────────────────────────────────────────────────────────────────
\section{Model 1 --- Core subscription ladder}

The spine. Five rungs instead of the current four, adding a downmarket
\plan{starter\_annual} so small contractors are not forced to choose between
``\Rs 2.4L'' and ``nothing''.

\begin{center}
\begin{longtable}{@{}L{2.0cm} L{3.9cm} R{2.3cm} L{1.6cm} L{3.4cm}@{}}
  \toprule
  \textbf{Tier} & \textbf{\texttt{plan\_id}} & \textbf{Annual (INR)} & \textbf{Active projects} & \textbf{Headline gating} \\
  \midrule
  \endhead
  Free trial   & \plan{trial}            & \Rs 0 / 14 days    & 1         & All templates, mobile, no AI \\
  Starter      & \plan{starter\_annual}  & \Rs 60{,}000       & 3         & Mobile + reports; no AI, no SSO \\
  Team         & \plan{team\_annual}     & \Rs 2{,}40{,}000   & unlimited & + AI reports, audit-log UI, SLA 1\,d \\
  Business     & \plan{business\_annual} & \Rs 6{,}00{,}000   & unlimited & + custom templates, SSO, 4\,h SLA \\
  Enterprise   & \plan{enterprise\_custom}& \Rs 12{,}00{,}000+ & unlimited & + CSM, on-prem, custom RBAC \\
  \bottomrule
\end{longtable}
\end{center}

\noindent\textbf{Unchanged rules from \texttt{PRICING.md}:} priced per workspace,
not per seat; annual upfront is default; multi-tenancy, mobile, soft-delete and
export are \emph{never} gated. \plan{starter\_annual} deliberately excludes AI and
SSO so it cannot cannibalise Team.

\subsection{Billing cadence as a price lever}
Cadence is sold, not given. Each step toward upfront cash earns the customer a
discount and earns you working capital.

\begin{center}
\begin{tabular}{@{}L{4.5cm} L{4.0cm} L{6.0cm}@{}}
  \toprule
  \textbf{Cadence} & \textbf{Adjustment} & \textbf{Notes} \\
  \midrule
  Annual upfront      & list price        & Default. Best margin + float. \\
  Quarterly           & \textbf{+5\%}     & Paid concession for cash-flow-shy buyers. \\
  Monthly bridge      & +0\%, 90 days max & Trial$\rightarrow$annual ramp only, never permanent. \\
  2-year prepaid      & \textbf{$-$20\%}  & Locks logo + reduces churn risk. \\
  3-year prepaid      & \textbf{$-$30\%}  & Floor; never discount deeper without a concession. \\
  \bottomrule
\end{tabular}
\end{center}

% ── Model 2 ──────────────────────────────────────────────────────────────────
\section{Model 2 --- Hybrid metered add-ons}

Flat platform fee \emph{plus} a meter on the few resources whose cost scales with
heavy use. This is the single biggest change to the revenue shape: it turns
power users into more revenue without raising the headline price that wins deals.

\begin{center}
\begin{tabular}{@{}L{4.0cm} L{3.2cm} R{2.6cm} L{4.6cm}@{}}
  \toprule
  \textbf{Metered SKU} & \textbf{\texttt{plan\_id} suffix} & \textbf{Unit price} & \textbf{Fair-use included} \\
  \midrule
  AI handover reports   & \plan{ai\_report}    & \Rs 500 / report  & 50/mo Team, unltd Business+ \\
  Document retention    & \plan{storage\_gb}   & \Rs 100 / GB-yr   & 25 GB Team, 100 GB Business \\
  SMS/WhatsApp alerts   & \plan{comms\_pack}   & \Rs 1 / message   & email always free \\
  Extra workspace       & \plan{extra\_ws}     & \Rs 90{,}000 / yr  & 1 per contract \\
  \bottomrule
\end{tabular}
\end{center}

\textbf{Mechanics.} Meters are counted against the \texttt{rate\_limits} table
that already backs \texttt{rate\_limit\_check}. AI reports are a natural meter ---
each generation is already concurrency-locked and rate-limited (10/hr). At
month close, sum overages above fair-use and raise a Razorpay invoice line.
Crucially, \textbf{fair-use is generous} so the meter is invisible to 90\% of
customers and only bites genuine power users --- the ones who can afford it.

\begin{keybox}
\textbf{Guardrail.} Never meter a \emph{measurement}. Metering test records or
equipment instances makes engineers under-report to dodge the meter --- it
corrupts the data you exist to protect. Meter only outputs (reports, storage,
notifications) and capacity (workspaces), never field input.
\end{keybox}

% ── Model 3 ──────────────────────────────────────────────────────────────────
\section{Model 3 --- Consumption / project packs}

For the seasonal contractor who refuses an annual commitment: sell prepaid
\textbf{project credits}. One credit unlocks one project workspace (full Team
features) for the life of that commissioning job. Credits never expire.

\begin{center}
\begin{tabular}{@{}L{2.6cm} L{2.8cm} R{1.8cm} R{2.6cm} L{2.6cm}@{}}
  \toprule
  \textbf{Pack} & \textbf{\texttt{plan\_id}} & \textbf{Credits} & \textbf{Price (INR)} & \textbf{Per project} \\
  \midrule
  Single  & \plan{pack\_1}  & 1  & \Rs 35{,}000   & \Rs 35{,}000 \\
  Crew    & \plan{pack\_5}  & 5  & \Rs 1{,}50{,}000 & \Rs 30{,}000 \\
  Fleet   & \plan{pack\_20} & 20 & \Rs 5{,}00{,}000 & \Rs 25{,}000 \\
  \bottomrule
\end{tabular}
\end{center}

\noindent\textbf{Why it works.} It converts the ``I'll wait until next month to
start'' objection into a purchase today, and it is a natural \emph{on-ramp to
annual}: a contractor who burns 8+ credits a year is shown the math --- ``you have
spent \Rs 2.0L on packs; \plan{team\_annual} is \Rs 2.4L for unlimited.'' The pack
price per project is deliberately \emph{higher} than the annual-implied rate so
the upgrade always looks rational.

\textbf{Schema fit.} Credits live as a balance in \texttt{companies.features}
(e.g. \flag{project\_credits: 5}); generating a project decrements it. No new
table required for v1.

% ── Model 4 ──────────────────────────────────────────────────────────────────
\section{Model 4 --- Outcome / value pricing (selective)}

Tie a slice of price to the value event itself: a \textbf{signed, approved
handover}. Offered only where a buyer resists fixed fees and wants
``pay-for-results''. Two shapes:

\begin{itemize}[noitemsep, leftmargin=1.4em]
  \item \textbf{Per commissioned bay} --- \Rs 1{,}500 per equipment instance that
        reaches \texttt{APPROVED}. Predictable for the buyer, scales with their
        real workload. Cap it at the Team annual so it can never exceed the flat
        alternative --- the cap is the upsell.
  \item \textbf{Per handover report delivered} --- \Rs 2{,}000 per
        \texttt{CLOSED} project's final report. Aligns spend with the deliverable
        the client actually pays \emph{them} for.
\end{itemize}

\noindent Use sparingly. Outcome pricing maximises trust on the first deal but
makes revenue lumpy and harder to forecast; migrate these accounts to flat annual
once they trust the product (typically renewal 1).

% ── Model 5 ──────────────────────────────────────────────────────────────────
\section{Model 5 --- Channel \& white-label (B2B2B)}

Let EPC majors and switchgear OEMs resell TestFlow to \emph{their}
subcontractors under their own brand. You sell wholesale; they own the
relationship.

\begin{center}
\begin{tabular}{@{}L{4.3cm} L{4.0cm} L{6.0cm}@{}}
  \toprule
  \textbf{Lever} & \textbf{Terms} & \textbf{Rationale} \\
  \midrule
  Wholesale rate     & list $\times$ \textbf{0.6}  & Partner keeps the margin + does the selling. \\
  White-label add-on & \Rs 1{,}00{,}000 / yr        & Custom domain, logo, colours (\flag{white\_label}). \\
  Minimum commit     & 10 workspaces / yr          & Filters out tyre-kickers; guarantees volume. \\
  Co-marketing       & case study + logo rights    & Lowers your CAC on the next direct deal. \\
  \bottomrule
\end{tabular}
\end{center}

\noindent\textbf{Schema fit.} Each reseller is a parent \texttt{company} record;
sub-tenants carry \plan{partner\_wholesale} as \texttt{plan\_id} and inherit the
\flag{white\_label} feature flag. Razorpay bills the \emph{partner} a single
consolidated invoice, not each sub-tenant. This is the highest-leverage channel
because one signed OEM can deliver 20--50 workspaces with near-zero direct CAC.

% ── Model 6 ──────────────────────────────────────────────────────────────────
\section{Model 6 --- Utility framework licensing}

The top of the market: PowerGrid, NTPC, state transcos. They do not buy
``a workspace'' --- they buy a \textbf{framework agreement} covering many sites for
several years, often with on-prem or private-cloud requirements.

\begin{itemize}[noitemsep, leftmargin=1.4em]
  \item \textbf{Site-block license} --- published Enterprise rate $\times$ 0.7 for
        20+ sites bought together (matches \texttt{SALES\_TARGETS} volume play).
  \item \textbf{Multi-year} --- 3--5 year terms, annual escalation clause (CPI or
        flat 5\%/yr) written in.
  \item \textbf{On-prem / private cloud} --- premium of \Rs 5{,}00{,}000+/yr for a
        dedicated Supabase project or self-hosted deployment; gated by the
        \flag{on\_prem} flag and a separate deployment SKU.
  \item \textbf{Contractual SLA \& security review} --- priced into the deal, not
        given away (each costs real founder/engineer time).
\end{itemize}

\noindent These are slow (6--12 month sales cycles) but anchor ARR and create
reference logos that close everything below them.

% ── Model 7 ──────────────────────────────────────────────────────────────────
\section{Model 7 --- Services \& data products}

Non-subscription revenue that improves margin and stickiness.

\begin{center}
\begin{tabular}{@{}L{5.2cm} R{3.2cm} L{5.5cm}@{}}
  \toprule
  \textbf{Item} & \textbf{Price (INR)} & \textbf{Type} \\
  \midrule
  Onboarding / setup fee        & \Rs 25{,}000 one-time  & Team \& Business; waivable as ``discount'' \\
  Data migration / import       & \Rs 25{,}000 one-time  & Legacy Excel $\rightarrow$ TestFlow \\
  Custom report template        & \Rs 50{,}000 each      & Per template (\flag{custom\_templates}) \\
  On-site training (8\,h)       & \Rs 50{,}000 + travel  & Per session \\
  Premium support (4\,h resp.)  & \Rs 2{,}00{,}000 / yr  & Upsell below Business SLA \\
  \textit{Benchmark / analytics pack} & \textit{future}  & Anonymised, opt-in failure/throughput data \\
  \bottomrule
\end{tabular}
\end{center}

\noindent The \emph{data product} is the long-game option: with enough tenants,
anonymised, opt-in commissioning benchmarks (mean time-to-handover, failure rates
by equipment type) become a sellable industry report --- a revenue stream with
\emph{zero} marginal delivery cost. Strictly opt-in; multi-tenancy isolation
(RLS) is never compromised for it.

% ── Packaging matrix ─────────────────────────────────────────────────────────
\section{Feature-gating matrix (maps to \texttt{companies.features})}

Every gate below is a boolean key in the \texttt{companies.features} JSONB,
enforced by \texttt{has\_feature()} server-side and \texttt{useFeature()} in the
UI. Flags default \texttt{TRUE} unless a downgrade sets them false --- so a missing
flag never accidentally locks a paying customer out.

\begin{center}
\begin{tabular}{@{}L{4.6cm} C{1.8cm} C{1.8cm} C{1.9cm} C{2.2cm}@{}}
  \toprule
  \textbf{Feature flag} & \textbf{Starter} & \textbf{Team} & \textbf{Business} & \textbf{Enterprise} \\
  \midrule
  \flag{ai\_reports}        & \no  & \yes & \yes & \yes \\
  \flag{audit\_log\_ui}     & \no  & \yes & \yes & \yes \\
  \flag{custom\_templates}  & \no  & \no  & \yes & \yes \\
  \flag{sso}                & \no  & \no  & \yes & \yes \\
  \flag{white\_label}       & \no  & add-on & add-on & \yes \\
  \flag{priority\_support}  & \no  & \no  & \yes & \yes \\
  \flag{advanced\_analytics}& \no  & \no  & \no  & \yes \\
  \flag{on\_prem}           & \no  & \no  & \no  & \yes \\
  \bottomrule
\end{tabular}
\end{center}

\noindent\textbf{Never gated (in every tier):} all test templates, mobile app,
multi-tenant RLS isolation, real-time / 30\,s polling sync, soft-delete \& restore,
PDF + Excel export. Audit-log \emph{capture} is always on (triggers); only the
\emph{viewer UI} is gated.

% ── Expansion / NRR ──────────────────────────────────────────────────────────
\section{Expansion levers \& the metric that matters}

A flat annual fee alone gives \tild{}100\% net revenue retention at best (you
only ever lose, never grow, within an account). The layered models above create
\textbf{negative churn} --- accounts that grow without re-negotiating the base:

\begin{enumerate}[noitemsep, leftmargin=1.6em]
  \item \textbf{Tier upgrades} --- Starter $\rightarrow$ Team when project cap bites.
  \item \textbf{Metered overages} --- AI/storage creep as usage deepens.
  \item \textbf{Extra workspaces} --- new business unit / region.
  \item \textbf{Pack$\rightarrow$annual conversion} --- the consumption on-ramp.
  \item \textbf{Add-ons} --- white-label, templates, premium support.
\end{enumerate}

\begin{keybox}
\textbf{The single number to track: ARR per workspace} (from \texttt{PRICING.md}),
now paired with \textbf{Net Revenue Retention}. If NRR $>$ 100\%, the expansion
models are working and you can grow ARR \emph{without} adding logos. Target NRR
$\geq$ 110\% by end of year 2.
\end{keybox}

% ── Implementation ───────────────────────────────────────────────────────────
\section{Implementation notes (what the code already supports)}

\begin{center}
\begin{tabular}{@{}L{5.0cm} L{9.5cm}@{}}
  \toprule
  \textbf{Need} & \textbf{Already in schema / what to add} \\
  \midrule
  Plan identity        & \texttt{subscriptions.plan\_id} (TEXT) --- use the slugs above. \\
  Feature gating       & \texttt{companies.features} JSONB + \texttt{has\_feature()} / \texttt{useFeature()}. \\
  Trial                & \texttt{companies.trial\_ends\_at} (14-day default trigger). \\
  Metering             & \texttt{rate\_limits} table + \texttt{rate\_limit\_check} RPC; sum overages monthly. \\
  Billing sync         & \texttt{razorpay-webhook} $\rightarrow$ \texttt{upsert\_subscription()} (idempotent). \\
  Seat count (optional)& \texttt{subscriptions.seat\_count} exists but stays 0 --- we do \emph{not} price per seat. \\
  \midrule
  \textbf{To build}    & project-credit balance in \texttt{features}; overage invoicing job; Razorpay plan objects per \texttt{plan\_id}; partner/parent-company billing roll-up. \\
  \bottomrule
\end{tabular}
\end{center}

% ── Anti-patterns ────────────────────────────────────────────────────────────
\section{Anti-patterns (carry-over + new)}

\begin{itemize}[noitemsep, leftmargin=1.4em]
  \item Never meter test records or equipment instances --- corrupts field data.
  \item Never gate multi-tenancy, mobile, or audit \emph{capture} --- core trust.
  \item Don't ship a permanent free tier; the trial + free template library is the lure.
  \item Don't let \plan{starter\_annual} include AI/SSO --- it would cannibalise Team.
  \item Don't price packs \emph{below} the annual-implied rate --- the upgrade must always look rational.
  \item Don't tag Enterprise ``most popular'' --- Team is the workhorse.
  \item Don't run outcome pricing forever --- migrate to flat annual by renewal 1.
\end{itemize}

\vspace{1em}
\hrule
\vspace{0.5em}
\noindent\textcolor{mgrey}{\small Companion docs: \texttt{PRICING.md} (strategy \&
display), \texttt{SALES\_TARGETS.md} (named accounts), financial-plan \texttt{.tex}
(unit economics \& 36-month ARR). Review quarterly; update slugs here when Razorpay
plan objects are created.}

\end{document}
